\chapter{DISCUSSION}

This study provides the first nationally representative quantification of socioeconomic inequality in HPV vaccine awareness and uptake among young women aged 15--29 in Vietnam. By integrating survey-weighted modeling with the Erreygers decomposition framework, we elucidate the structural determinants that drive the profound equity gap in cervical cancer prevention.

\textbf{The knowing--doing chasm.}
Our analysis reveals a striking dichotomy between cognition and action: while 62.4\% of women have heard of the HPV vaccine, a mere 7.5\% have received it. This 55-percentage-point awareness--uptake gap underscores a fundamental failure of translation from health literacy to health behavior. Similar implementation bottlenecks have been documented across Southeast Asia---Thailand reports 58.3\% awareness against sub-5\% uptake \cite{songthap2021}, and the Philippines exhibits single-digit coverage despite 55.7\% awareness \cite{gabuyo2022}. In stark contrast, high-income nations with publicly financed, school-based programs routinely achieve coverage exceeding 70\% \cite{marlow2019}. The Vietnamese gap is therefore not merely informational but fundamentally structural, rooted in the absence of universal public financing and the prohibitive out-of-pocket costs of the vaccine series.

\textbf{The anatomy of inequality.}
The concentration indices demonstrate severe pro-rich inequality, with the magnitude of disparity for vaccination (CI = 0.389) more than double that for awareness (CI = 0.149). This differential implies that while information diffuses across socioeconomic strata---albeit imperfectly---the actual receipt of the vaccine is profoundly constrained by barriers at the point of service delivery. This aligns with Penchansky and Thomas's access framework \cite{penchansky1981}. Information can diffuse relatively inexpensively through digital networks; receiving a vaccine, however, demands physical proximity to a health facility, the financial capacity to absorb high costs, and the operational autonomy to seek care---dimensions sharply patterned by socioeconomic position.

Critically, the negative CI for the awareness--vaccination gap ($-$0.035) reinforces this paradigm: among women already possessing knowledge of the vaccine, the poorest are disproportionately paralyzed from acting upon it. Policy interventions predicated solely on raising awareness will therefore inevitably fail to close the equity divide.

\textbf{Wealth as the dominant engine of disparity.}
The Erreygers decomposition unequivocally isolates household wealth as the primary catalyst of inequity, explaining 38.8\% of awareness inequality and an overwhelming 51.4\% of vaccination inequality. This dominance mirrors global multi-country analyses of childhood immunization \cite{hosseinpoor2012}. In the Vietnamese context, where the HPV vaccine series costs approximately 1.5--3.0 million VND in the private sector \cite{nguyen2021}, the intervention effectively functions as a luxury good rather than a basic public health right. Adjusted multivariable models confirm this severe penalty: women in the poorest quintile face an 88\% reduction in the probability of vaccination (PR = 0.118) compared to their wealthiest counterparts.

\textbf{Education: the second pillar of inequality.}
Education emerged as an important driver, contributing 7.7\% and 12.3\% to awareness and vaccination inequality, respectively. The dose-response gradient is steep: primary-educated women exhibit a 53\% reduction in awareness likelihood (PR = 0.475) and lower secondary women show a 58\% reduction in vaccination (PR = 0.422) relative to those with upper-secondary education or higher. These associations likely operate through enhanced health literacy and stronger integration into social networks that normalize preventive health behaviors \cite{nutbeam2008}.

\textbf{The role of the information ecosystem and geographic fault lines.}
Mass media exposure contributed 11.1\% and 23.8\% to the inequality in awareness and vaccination, respectively. Women lacking exposure to standard media channels---who are disproportionately impoverished, rural, and ethnically minoritized---are effectively excluded from the primary circuits of health communication. Ethnic minority women report the lowest awareness (36.1\%) and near-zero vaccination (1.6\%), reflecting the deep marginalization of Vietnam's remote populations \cite{pham2020}. This divide necessitates community-based outreach leveraging local health workers and peer networks \cite{wakefield2010}.

\section{Strengths and limitations}

This study draws on a large (N~=~4,102), nationally representative survey with a well-characterized sampling design, enabling robust population-level inference for young Vietnamese women. The use of survey-weighted Poisson regression yields directly interpretable prevalence ratios, circumventing the overestimation of risk inherent in odds ratios for prevalent outcomes like awareness \cite{barros2003}. The Erreygers decomposition framework provides precise, policy-actionable quantification of the specific drivers of inequality.

However, several limitations must be acknowledged. First, the cross-sectional design precludes causal inference. Second, outcome measures are self-reported and may be susceptible to recall or social desirability bias. Third, the MICS6 dataset lacks granular clinical details such as the number of doses received, timing relative to sexual debut, or reasons for non-vaccination. Fourth, by restricting the sample to women aged 15--29, the findings specifically characterize inequality among young women and may not generalize to older cohorts. Finally, ICT exposure was not excluded from models but its near-universal prevalence (97.9\%) limits the ability to detect heterogeneity in digital information exposure. Future longitudinal research incorporating qualitative methodologies is required to decode the behavioral nuances of vaccine hesitancy and structural exclusion in Vietnam.